\paragraph{A rational certificate for the transition value.}
For $h=2\pi/\sqrt{47}$ put $G(a)=\int_1^\infty e^{-av^2}\,dv$.
The substitution $v=e^{u/2}$ in the centered Mellin integral gives
$\Lambda_j(1/2)=4\sum_{n\ge1}a_j(n)G(hn)$.
For $v\ge1$, $v^2\ge1+2(v-1)$, hence
$G(hn)\le e^{-hn}/(2hn)$.
Using $|a_j(n)|\le n$ and $e^{-h}<r=41/100$, the absolute tail after $N=28$,
before multiplication by $4$, is bounded by
$r^{N+1}/(2h(1-r))$. The finite coefficients come from enumerating the three
binary quadratic forms in the theta identities.

For each $a>0$ the program subtracts
$\int_0^1e^{-av^2}\,dv$ from $\tfrac12\sqrt{\pi/a}$.
The integrated Taylor polynomials
$\sum_{k=0}^{m}(-a)^k/[k!(2k+1)]$ are an upper bound at even $m=122$ and a
lower bound at $m=123$, by the sign of the Taylor remainder of $e^{-x}$ for
$x\ge0$. The program evaluates the two endpoints of the interval for $a$
using monotonicity of $G$. It bounds $\pi$ by Machin's arctangent identity,
and all square roots by integer square roots after rational scaling. The
enclosures from \texttt{check\_threshold.py} are
\[
\begin{split}
 0.6660152520&\le\Lambda_1(1/2)\le0.6660152521,\\
 0.4586804524&\le\Lambda_2(1/2)\le0.4586804526,\\
 1.4520244940&\le C_0\le1.4520244943 .
\end{split}
\]
The last line uses the unrounded intervals. For example, the pencil with
$c=7/5$ has two simple real zeros, while the coefficient $3/2$
gives none.
